\subsection{Corpus, Coding and How to Read the Numbers}
\label{sec:empirical-corpus}

The corpus is the SMOTE25 collection selected in Section~\ref{sec:methodology}: 76 papers after duplicate copies were merged. Of these, 71 report an empirical evaluation on imbalanced data; the other five are two surveys~\parencite{pradipta2021smote,hassanat2022stop} and three foundational works that predate imbalance-aware evaluation~\parencite{wilson1972asymptotic,tomek1976modifications,freund1997adaboost}. Among the 71 empirical studies, 39 evaluate binary problems, 19 multi-class problems and 13 both (Figure~\ref{fig:empirical-corpus}). Thirty-seven of the 71 are SMOTE-family methods in the sense of Section~\ref{sec:taxonomy} (34 core and 3 partial members: 24 binary only, 5 multi-class only and 8 for both settings); the others are competing resamplers, ensembles, cost-sensitive methods, generative models and empirical studies that SMOTE-family papers compare with. We report results for all 71 and point out where the 37 family members behave differently.

\begin{figure}[t]
  \centering
  \includegraphics[width=\linewidth]{empirical-corpus}
  \caption{Empirical Studies in the Corpus by Publication Period and Problem Setting}
  \Description{Stacked columns for seven publication periods from 1997 to 2026; each column splits the empirical studies of that period into binary, multi-class and both.}
  \label{fig:empirical-corpus}
\end{figure}

To show change over time we group the studies into four periods of similar size: 1997 to 2009 (14 studies), 2010 to 2015 (13), 2016 to 2020 (20) and 2021 to 2026 (24). A statement such as ``G-mean rose from 29\% to 71\%'' means that 4 of the 14 studies of the first period and 17 of the 24 studies of the last period report G-mean. With periods this small, one study moves a share by 4 to 8 points, so we only call a change a trend when it exceeds 15 points or when a practice appears or disappears completely.

For each study we coded, from the full text and supplementary files, the benchmark data, the base learners, the competing methods, every measure together with its exact definition, the validation protocol, how hyperparameters were chosen, the statistical tests and the form in which results are reported. A field the paper does not state was coded as not reported, never inferred. A second reading recorded why each study made its choices: the reasons it gives for its data, learners and measures, the measure and the data used to choose its settings, the methodological sources it cites, and any evidence about the resampled data that does not come from a classifier. Every code from the second reading carries a page reference and a short verbatim quote, so it can be checked against the paper \TODO{a co-author to verify the second-reading quotes before submission}. We also recorded which software a study names (for example Weka, KEEL, MATLAB, scikit-learn or imbalanced-learn). Only named software is counted, so these counts are lower bounds. Table~\ref{tab:empirical-studies} in the appendix lists the main codes of every study.

When we explain why a practice changed, we use three levels of evidence and say which one applies. Direct evidence means that the studies themselves give the reason, or that the link can be measured in the corpus, for example the share of composite-measure users who still report recall separately. Circumstantial evidence means that timing and association point the same way but no study states the link, for example the parallel rise of Python and of CART. A conjecture is a mechanism we propose that the corpus cannot confirm. None of these is a causal estimate in the statistical sense; with 71 studies and no intervention, the corpus can show that a change and its proposed cause move together and that the studies say so, not more.

To see whether careful evaluation is rewarded, we took citation counts from OpenAlex on 18 October 2026, available for 70 of the 71 studies, and compare the median number of citations of the studies that follow a practice with that of the studies that do not (Section~\ref{sec:empirical-why}). We report the raw counts. A count has to be read against the year of publication, because older papers have had more time to be cited: the median study published up to 2009 has 1,738 citations, against 21 for the studies published since 2021. Every comparison therefore gives the median year of each group next to its median count, and where the two groups differ in age we also compare studies of the same period.

\paragraph{Relation to the fifteen-year review.} This analysis complements the review that marked the fifteenth anniversary of SMOTE~\parencite{fernandez2018smote}. That review catalogued more than 85 extensions by seven properties of their design (how the instances to oversample are selected, integration with undersampling, the type of interpolation, operation under changes of dimensionality, adaptive generation, relabelling and the filtering of noisy synthetic instances), listed SMOTE-based ensembles and 22 SMOTE-based methods for 12 other learning paradigms, and set out the open challenges: small disjuncts, noise, class overlap, dataset shift, dimensionality, real-time processing and big data. On evaluation it says three things. It recalls why accuracy failed on the data that motivated SMOTE, where a decision tree reached about 97\% accuracy on the mammography data while always predicting the majority class gave 97.68\%, and the two kinds of error had different costs, so the trade-off between true and false positives had to be shown with ROC analysis. It summarises the empirical comparisons available in 2018. And it recommends a validation scheme that is robust to the dataset shift induced by partitioning. It does not examine which measures, protocols or tests the studies use, which is what this section adds. Table~\ref{tab:empirical-review15} lists the points of the review that bear on evaluation and what the studies published since did; we return to each of them below (Sections~\ref{sec:empirical-baselines}, \ref{sec:empirical-direct}, \ref{sec:empirical-protocol} and~\ref{sec:empirical-why}).

\input{tables/empirical-review15}

